\section{An unconditional genus-three instance: the Fermat quartic}
\label{sec:fermat-instance}

The general finiteness theorem remains unresolved in the requested setting.
There is nevertheless a substantial fixed-curve instance for which the
entire arithmetic argument can be given by primitive recursive operations
and quantifier-free induction.  We prove
\[
 C:\ X^4+Y^4=Z^4,
 \qquad
 C(\mathbb Q)=
 \{[1:0:1],[-1:0:1],[0:1:1],[0:-1:1]\}.
 \tag{FQ}\label{eq:fq-points}
\]
This is a classification, with a height bound and an explicit complete list.
It is stronger than merely ruling out arbitrarily long lists of distinct
points.  The arithmetic proof does not invoke Faltings' theorem.

The elementary descent below follows the mathematical route also recorded
in the inspected \texttt{diophantine\_geometry\_ipra.zip} draft, but every
ingredient needed here is proved.  The argument is a human-readable
arithmetic proof; this section does not claim that its deeply embedded
\texttt{Prf PRAAx} proof object has already been constructed.  The ordinary
geometric identification of this particular curve as a smooth genus-three
curve is separated at the end.

\subsection{Finite arithmetic conventions}

Signed integers are represented by a sign and a magnitude, with a unique
code for zero.  A rational number is represented by a reduced pair
$(u,v)$ with $u\in\mathbb Z$, $v\in\mathbb N$, $v>0$ and
$\gcd(|u|,v)=1$.  A rational projective point is a nonzero rational triple
modulo multiplication by a nonzero rational.  Its canonical code is an
integer triple $(X,Y,Z)$ with
\[
 \gcd(|X|,|Y|,|Z|)=1
 \quad\hbox{and first nonzero coordinate positive}.
 \tag{F1}
\]
Here the three-variable gcd means the gcd of the first two magnitudes
followed by the third.  Multiply the rational coordinates by the product
of their positive denominators, divide the resulting integers by their
positive gcd, and fix the sign.  This produces (F1) by a specified finite
algorithm.  Homogeneity preserves the quartic equation through these
operations.

For completeness, this representative is unique.  Suppose two primitive
integer triples differ by a rational factor $u/v$, with $v>0$ and
$\gcd(|u|,v)=1$.  From $vP_i=uQ_i$, Euclid's divisibility lemma, proved
below, implies that $v$ divides every $Q_i$.  Primitivity gives $v=1$.
The same argument in the other direction gives $|u|=1$, and the sign
convention gives $u=1$.  Every assertion here reduces to integer arithmetic
on three entries.  Define the nonlogarithmic height of a canonical point by
\[
 H(X,Y,Z)=\max(|X|,|Y|,|Z|).
\]

The calculations use addition, multiplication, comparison, parity, bounded
division, bounded gcd, and bounded square-root searches.  All are totalized
on malformed inputs, for example by returning zero.  A division $a/b$ in a
valid branch below is used only after $b>0$ and $b\mid a$ have been proved.
It can be implemented by testing $q\leq a$ for $bq=a$.  An exact square
root bounded by $N$ is obtained by testing the integers $r\leq N$ for
$r^2=a$.  Thus no unbounded minimization is concealed in the notation.

\subsection{The gcd and square-splitting lemmas}

The zero branches needed for projective normalization are fixed explicitly:
$\gcd(a,0)=a$, $\gcd(0,b)=b$, and $\gcd(0,0)=0$ for nonnegative
inputs. When one input is positive and the other is zero, the Bezout
coefficient pairs are respectively $(1,0)$ and $(0,1)$, and the
common-divisor properties follow immediately. A nonzero triple therefore
has a positive three-variable gcd even when two coordinates are zero.
The remaining computation is the positive-input case proved next.

\begin{lemma}[Bounded gcd and its certificates]\label{lem:fq-gcd}
For positive integers $a,b$ one can compute $g>0$ and signed integers
$u,v$ such that
\[
 g\mid a,\qquad g\mid b,\qquad ua+vb=g,
\]
and every common divisor of $a,b$ divides $g$.  These assertions have
primitive recursive witnesses and proofs by quantifier-free induction.
\end{lemma}

\begin{proof}
Start with $(a_0,b_0)=(\max(a,b),\min(a,b))$.  If $b_i>0$, perform bounded
division $a_i=q_i b_i+r_i$, $0\leq r_i<b_i$, and set
$(a_{i+1},b_{i+1})=(b_i,r_i)$.  If $b_i=0$, leave the pair fixed.
Before it stops, the positive sum $a_i+b_i$ strictly decreases: indeed
$b_i+r_i<a_i+b_i$.  Hence $a+b+1$ iterations suffice to reach a zero
second entry.  The first entry is positive throughout the stopped state.

Track both entries as signed linear combinations of the original inputs.
At a division step, replace the coefficient pair for the first entry by
that for the second, and replace the coefficient pair for the second by
the old first pair minus $q_i$ times the old second pair.  These rules
preserve the two displayed linear-combination equations.  If
$M=\max(a,b)$, then $q_i\leq M$, and every tracked coefficient after $j$
steps has magnitude at most $(M+1)^{j+1}$.  In particular, a finite storage
bound is available, although primitive recursion does not need one for
its definition.

For any fixed divisor $d$, the equations
$a_i=q_i b_i+r_i$ and $r_i=a_i-q_i b_i$ show that
\[
 (d\mid a_i\ \wedge\ d\mid b_i)
 \ \longleftrightarrow\
 (d\mid b_i\ \wedge\ d\mid r_i).
\]
At the stopped pair $(g,0)$ this proves both that every original common
divisor divides $g$ and, on taking $d=g$, that $g$ divides both original
inputs.  The coefficient invariant supplies $u,v$.

All induction invariants used in this argument are quantifier-free once
divisibility is replaced by its bounded characteristic function.  The
divisor $d$ is a free numerical parameter during induction, not an
unbounded quantifier in its induction formula.  Termination can use the
quantifier-free invariant
$b_i=0\ \vee\ (a_i+b_i+i\leq a+b)$ and the fixed fuel $a+b+1$.
\end{proof}

If $g=\gcd(a,b)$, the positive integers $a/g$ and $b/g$ are coprime:
a common divisor $d$ would make $gd$ a common divisor of $a,b$, hence
$gd\mid g$ and $d=1$.  This supplies the gcd normalization used below.

\begin{lemma}[Coprime cancellation and square extraction]
\label{lem:fq-squares}
The following implications have primitive recursive witness operations.
\begin{enumerate}
\item If $\gcd(a,b)=1$ and $a\mid bc$, then $a\mid c$.
\item If $d>0$ and $d^2\mid c^2$, then $d\mid c$.
\item If $a,b,c>0$, $\gcd(a,b)=1$, and $ab=c^2$, then there are
positive integers $r,s$ with $a=r^2$, $b=s^2$, and $c=rs$.
\end{enumerate}
\end{lemma}

\begin{proof}
For (1), multiply a Bezout identity $ua+vb=1$ by $c$.
Both terms on the right of $c=uac+vbc$ are divisible by $a$.
We also need two consequences of Bezout identities.  Squaring
$ua+vb=1$ and collecting terms gives
\[
 1=a(u^2a+2uvb)+v^2b^2,
\]
so coprimality of $a,b$ implies coprimality of $a,b^2$.
Multiplying a Bezout identity for $a,b$ by one for $a,c$ shows that
coprimality of $a$ with each of $b,c$ implies coprimality of $a,bc$.

For (2), the case $c=0$ is immediate.  Otherwise set
$h=\gcd(c,d)$, $c=hr$, and $d=hs$.  The normalized pair $r,s$ is
coprime.  Cancelling $h^2$ in the divisibility assumption gives
$s^2\mid r^2$, hence $s\mid r^2$.  But $s$ is coprime to $r^2$ by
the preceding Bezout calculation.  Applying (1), or its same Bezout
identity, gives $s\mid1$, so $s=1$.  Therefore $d=h\mid c$.

For (3), put $h=\gcd(a,c)$, $a=hu$ and $c=hv$.
Then $u,v$ are positive and coprime, and cancellation in $ab=c^2$
gives $ub=hv^2$.  Since $u$ is coprime to $v^2$, (1) implies
$u\mid h$.  Write $h=ue$.  Substitution and cancellation now yield
\[
 a=eu^2,\qquad b=ev^2,\qquad c=euv.
\]
The positive integer $e$ divides both $a,b$, so $e=1$.  Taking
$r=u$ and $s=v$ proves the result.  The construction uses only gcd
and exact division.  In particular, it is a witness-producing proof,
without an appeal to the existence of prime factorizations.
\end{proof}

The last lemma can be applied repeatedly to a product of any fixed number
of pairwise coprime positive factors.  Each factor is coprime to the
product of the others by the Bezout product calculation, so a square
product of four such factors makes all four factors squares.  This is
three uses of the lemma, not an unbounded search for square factors.

\subsection{Normalization and primitive Pythagorean parameters}

Call $(a,b,c,s)$ a square-area triangle if all four entries are positive
integers and
\[
 a^2+b^2=c^2,\qquad ab=2s^2.
 \tag{F2}\label{eq:fq-area}
\]
It is primitive if $\gcd(a,b)=1$.

\begin{lemma}[Primitive normalization]\label{lem:fq-normalize}
Every square-area triangle gives a primitive square-area triangle
$(A,B,C,S)$ with $C\leq c$ by a primitive recursive construction.
\end{lemma}

\begin{proof}
Put $g=\gcd(a,b)$, $A=a/g$, and $B=b/g$.  Then
$g^2\mid c^2$ by the Pythagorean equation, so
Lemma~\ref{lem:fq-squares}(2) gives $g\mid c$.  Set $C=c/g$.
We have $A^2+B^2=C^2$ and $\gcd(A,B)=1$.
The legs cannot both be even.  They cannot both be odd either, since
squares modulo $4$ are $0$ or $1$, whereas the sum of two odd squares
is $2$ modulo $4$.  Exactly one of $A,B$ is therefore even.
Consequently $K=AB/2$ is an integer, and the second equation in
\eqref{eq:fq-area} becomes $s^2=g^2K$.  The same square-divisibility
lemma gives $g\mid s$.  Set $S=s/g$ and cancel $g^2$ in both equations.
All output entries are positive, and $C=c/g\leq c$.
\end{proof}

\begin{lemma}[Explicit Pythagorean parameters]\label{lem:fq-pyth}
If $A,B,C>0$, $A^2+B^2=C^2$, and $\gcd(A,B)=1$, exchange the legs
if necessary to make $A$ even.  There are positive integers $m>n$ with
\[
 A=2mn,\qquad B=m^2-n^2,\qquad C=m^2+n^2,
 \tag{F3}\label{eq:fq-pyth}
\]
such that $\gcd(m,n)=1$ and $m,n$ have opposite parity.
The parameters are computable by searches bounded by $C$.
\end{lemma}

\begin{proof}
The parity argument just given makes $B$ and $C$ odd.
Also $B,C$ are coprime.  For if $d$ divides them both, then
$d^2\mid C^2-B^2=A^2$, so $d\mid A$ by
Lemma~\ref{lem:fq-squares}(2), and $d=1$.
Since $A>0$, the equation gives $C>B>0$.  Thus
\[
 U=(C+B)/2,\qquad V=(C-B)/2
\]
are positive integers with $U>V$.  A common divisor of $U,V$
divides $U+V=C$ and $U-V=B$, hence equals $1$.  Moreover
\[
 UV=(C^2-B^2)/4=(A/2)^2.
\]
Lemma~\ref{lem:fq-squares}(3) gives $U=m^2$, $V=n^2$, and
$A/2=mn$.  These equations give \eqref{eq:fq-pyth} and $m>n$.
A common divisor of $m,n$ would divide $U,V$, so they are coprime.
Their squares sum to the odd number $C$, so their parities differ.
Because $U,V\leq C$, their nonnegative square roots are at most $C$.
They are unique: if $r<t$ are nonnegative, then
$t^2-r^2=(t-r)(t+r)>0$.  A scan of $0,\ldots,C$ therefore returns
the required roots and has a proved success condition.
\end{proof}

\subsection{A descending map, with a numerical stopping bound}

\begin{lemma}[One explicit descent step]\label{lem:fq-step}
A primitive square-area triangle $(A,B,C,S)$ gives another primitive
square-area triangle $(A',B',C',S')$ with $0<C'<C$.
The step is a total primitive recursive function after assigning zero
output to invalid inputs.
\end{lemma}

\begin{proof}
Use the ordering and parameters of Lemma~\ref{lem:fq-pyth}.
Substituting \eqref{eq:fq-pyth} in $AB=2S^2$ and cancelling $2$
gives
\[
 mn(m+n)(m-n)=S^2.
 \tag{F4}
\]
All four factors are positive and pairwise coprime.
For pairs involving $m$ or $n$, subtraction reduces a common divisor
to a common divisor of $m,n$.  For the remaining pair, a common divisor
of $m+n,m-n$ divides $2m,2n$, hence divides $2$ by a Bezout identity
for $m,n$.  Both $m+n,m-n$ are odd, so that divisor is $1$.
Square splitting supplies positive integers $r,s,t,u$ such that
\[
 m=r^2,\qquad n=s^2,\qquad m+n=t^2,\qquad m-n=u^2.
 \tag{F5}\label{eq:fq-four-squares}
\]
The integers $t,u$ are odd, $t>u>0$, and they are coprime, since
any common divisor would divide the coprime pair $m+n,m-n$.
The integer $s$ is even.  Indeed, if $m$ were even and $n$ odd,
then $m=r^2\equiv0\pmod4$ and $n=s^2\equiv1\pmod4$,
making $m-n=u^2\equiv3\pmod4$, impossible.  Hence $n$ is even,
and so its square root $s$ is even.

Define
\[
 A'=(t+u)/2,\qquad B'=(t-u)/2,\qquad
 C'=r,\qquad S'=s/2.
 \tag{F6}\label{eq:fq-step}
\]
These are positive integers.  Direct expansion gives
\[
 (A')^2+(B')^2=(t^2+u^2)/2=m=r^2=(C')^2,
\]
and
\[
 A'B'=(t^2-u^2)/4=n/2=s^2/2=2(S')^2.
\]
A common divisor of $A',B'$ divides their sum $t$ and their
difference $u$, so the new triangle is primitive.  Finally,
\[
 0<C'=r\leq r^2=m<m^2+n^2=C.
 \tag{F7}\label{eq:fq-decrease}
\]
The last inequality follows from $m\geq1$, $n\geq1$ and
$m^2\geq m$.

For a completely specified algorithm, first test primitivity and both
equations.  On failure return $(0,0,0,0)$.  On success, exchange the
legs to make the first even, find $m,n$ as the unique exact roots of
$(C+B)/2,(C-B)/2$, and find $r,s,t,u$ as the exact roots in
\eqref{eq:fq-four-squares}.  Return \eqref{eq:fq-step}.
Every root search can use the input bound $C$: both $m,n\leq C$,
and $m+n\leq m^2+n^2=C$; the other radicands are no larger.
All denominators in the valid branch are the numeral $2$ and have
just been proved to divide their numerators.  Exact-root failure can
also be totalized by returning zero, and the proof shows that this
failure branch is never taken on a valid input.
\end{proof}

\begin{theorem}[No square-area triangle in intuitionistic PRA]
\label{thm:fq-no-area}
There are no positive integers satisfying \eqref{eq:fq-area}.
\end{theorem}

\begin{proof}
By Lemma~\ref{lem:fq-normalize}, a supposed solution produces a
primitive one, encoded by a state $q$.  Let $G(q)$ be the Boolean
predicate of being such a primitive solution, let $D$ be the total
step function from Lemma~\ref{lem:fq-step}, and let $\rho(q)$ be its
hypotenuse coordinate.  Define
\[
 q_0=q,\qquad q_{i+1}=D(q_i)
\]
by primitive recursion.  Under the assumption $G(q)$, induction gives
\[
 G(q_i)\quad\wedge\quad \rho(q_i)+i\leq\rho(q).
 \tag{F8}\label{eq:fq-invariant}
\]
The base is immediate.  For the step, Lemma~\ref{lem:fq-step} gives
$G(q_{i+1})$ and
$\rho(q_{i+1})+1\leq\rho(q_i)$; adding $i$ proves the inequality
at $i+1$.  This is quantifier-free induction: $G$ and the comparison
are fixed primitive recursive Boolean tests, and the state iteration
is a fixed primitive recursive function.

Now instantiate \eqref{eq:fq-invariant} at the specified numeral-valued
term $i=\rho(q)$.  Its inequality forces $\rho(q_i)=0$, whereas
$G(q_i)$ requires $\rho(q_i)>0$.  The contradiction discharges the
assumption $G(q)$.  The argument uses exactly the initial hypotenuse
as descent fuel and does not select a least counterexample or an
infinite sequence.  Negation is proved by contradiction from the
assumption being negated; no double-negation elimination is used.
\end{proof}

\subsection{From descent to the complete quartic list}

\begin{lemma}[Fourth-power difference]\label{lem:fq-difference}
There are no positive integers $x>y$ and $z$ with $x^4-y^4=z^2$.
Consequently, $a^4+b^4=c^4$ has no solution in positive integers.
\end{lemma}

\begin{proof}
A solution of $x^4-y^4=z^2$ would produce the positive integer
triangle
\[
 a=2x^2y^2,\qquad b=x^4-y^4,\qquad c=x^4+y^4,
 \qquad s=xyz.
\]
Expanding shows
\[
 (2x^2y^2)^2+(x^4-y^4)^2=(x^4+y^4)^2,
 \qquad ab=2x^2y^2z^2=2s^2,
\]
contradicting Theorem~\ref{thm:fq-no-area}.
If $a,b,c>0$ satisfy $a^4+b^4=c^4$, then $c>a$ and
$c^4-a^4=(b^2)^2$, a forbidden fourth-power difference.
\end{proof}

\begin{theorem}[The exact rational points of the Fermat quartic]
\label{thm:fq-complete}
In canonical integer projective codes, intuitionistic PRA proves that
every point of $X^4+Y^4=Z^4$ is one of
\[
 (1,0,1),\quad (1,0,-1),\quad (0,1,1),\quad (0,1,-1).
 \tag{F9}\label{eq:fq-canonical}
\]
Each listed triple is a point, the triples are pairwise distinct, and
every point has height $1$.  Therefore \eqref{eq:fq-points} holds.
\end{theorem}

\begin{proof}
Take a nonzero canonical integer triple satisfying the equation.
If $Z=0$, nonnegativity of fourth powers gives $X=Y=0$, a
contradiction.  Thus $|Z|>0$.  Equality of integers is decidable.
If both $X$ and $Y$ were nonzero, the positive integers
$|X|,|Y|,|Z|$ would contradict Lemma~\ref{lem:fq-difference}.
Case distinction on $X=0$ and $Y=0$ therefore leaves a zero
coordinate among $X,Y$.

If $Y=0$, the equation is $|X|^4=|Z|^4$.  Strict monotonicity of
the fourth power on nonnegative integers gives $|X|=|Z|$.
Primitivity forces this common positive value to be $1$.  The
first-nonzero-positive convention gives $X=1$, and $Z=1$ or $Z=-1$.
These are the first two triples of \eqref{eq:fq-canonical}.
If $X=0$, the identical argument gives $Y=1$ and $Z=1$ or $Z=-1$,
the last two triples.  They exhaust all cases.

Substitution verifies the equation for each listed triple, and direct
coordinate comparison proves their distinctness and their height $1$.
The second canonical triple represents $[-1:0:1]$, and the fourth
represents $[0:-1:1]$, by multiplication by $-1$.
All uses of case distinction concern explicit integer comparisons;
every induction above has a quantifier-free Boolean invariant.
\end{proof}

\subsection{The usual genus-three interpretation}

This paragraph identifies the arithmetic equation with a curve in the
ordinary statement of Faltings' theorem.  It is not an extra premise in
the integer classification just proved.

Over any characteristic-zero field, the three partial derivatives of
$F=X^4+Y^4-Z^4$ are $4X^3,4Y^3,-4Z^3$.  They cannot vanish
simultaneously at a projective point: their vanishing would force
$X=Y=Z=0$.  The projective Jacobian criterion therefore makes the
quartic smooth.  Absolute irreducibility can also be checked directly.
Over any extension field $E$, dehomogenize to
$f(X,Y)=Y^4+X^4-1$.  If $f=A B$ were a nontrivial factorization in
$E[X,Y]$, the leading coefficients in $Y$ would multiply to $1$.
After scaling, both factors would be monic in $Y$ and of positive
$Y$-degree; a factor of $Y$-degree zero would be a unit.
Reducing at $X=1$ gives $A(1,Y)B(1,Y)=Y^4$.  Both monic factors
of positive degree must have zero constant term, since the only
monic divisors of $Y^4$ are powers of $Y$.  Hence $(X-1)$ divides
both $A(X,0)$ and $B(X,0)$, so $(X-1)^2$ divides $X^4-1$.
This is impossible because
\[
 X^4-1=(X-1)(X^3+X^2+X+1),
 \qquad 1+1+1+1=4\ne0.
\]
A homogeneous factor which becomes a unit after setting $Z=1$
would be a power of $Z$ times a scalar, but $Z$ does not divide $F$.
Thus the homogeneous quartic is geometrically integral as well.

Its genus is $3$.  One can see the numerical calculation from the
Hilbert polynomial: for $n\geq4$, multiplication by $F$ is injective
on degree $n-4$ homogeneous polynomials, so the degree $n$ part of
$E[X,Y,Z]/(F)$ has dimension
\[
 \binom{n+2}{2}-\binom{n-2}{2}=4n-2=4n+1-3.
\]
For a smooth geometrically integral projective plane curve, the genus
is the number $g$ in this polynomial $dn+1-g$.  Equivalently, the
plane-curve genus formula gives $(4-1)(4-2)/2=3$; the usual geometric
theorem and its cohomological proof are recorded in the Stacks Project,
Section~53.9, Lemma~53.9.3, at
\url{https://stacks.math.columbia.edu/tag/0BYA}.
Thus \eqref{eq:fq-points} is a genuine genus-greater-than-one instance,
while its arithmetic proof requires none of this general geometry.

\subsection{Concrete Lean completion targets for this instance}

The proof gives a finite development with no missing finiteness theorem.
Implement the signed-number and tuple codes, and each fixed arithmetic
operation, as terms in the repository's \texttt{PR} syntax.  Prove
the bounded Euclidean and exact-root invariants, then the three
square lemmas, primitive normalization, the parameter equations,
and the preservation and decrease theorems for the displayed map $D$.
The required invariant for \texttt{PRA.qfInd} is exactly
\eqref{eq:fq-invariant}.  The final classification is a universal
formula whose matrix is a Boolean validity implication followed by
the four explicit coordinate equalities in \eqref{eq:fq-canonical}.

The desired final objects are an actual derivation
\[
 []\vdash[\mathtt{PRAAx}]\ \varphi_{\rm quartic}
\]
and its image under \texttt{pra\_iff\_native}.  The repository's
ambient theorem \path{LeanProofs.fermat_four_no_positive_nat_solutions}
is useful for comparison, but its existence does not construct that
object-theory derivation.  Conversely, completing this specific
arithmetic development would not complete the general Faltings
formalization: the explicit descent map is special to this quartic.
